% Generated draft: requires mathematical and source review.
\documentclass[11pt]{article}
\usepackage[margin=25mm]{geometry}
\usepackage{fontspec}
\setmainfont{Latin Modern Roman}
\usepackage{amsmath,amssymb,mathtools,unicode-math}
\setmathfont{Latin Modern Math}
\usepackage{enumitem,needspace,xurl,hyperref}
\setlength{\parindent}{0pt}
\setlength{\parskip}{5pt}
\setlength{\emergencystretch}{3em}
\newcommand{\N}{\mathbb{N}}
\newcommand{\Z}{\mathbb{Z}}
\newcommand{\Q}{\mathbb{Q}}
\newcommand{\R}{\mathbb{R}}
\newcommand{\C}{\mathbb{C}}
\newcommand{\F}{\mathbb{F}}
\newcommand{\A}{\mathbb{A}}
\newcommand{\PP}{\mathbb{P}}
\newcommand{\E}{\mathbb{E}}
\newcommand{\eps}{\varepsilon}
\newcommand{\dd}{\,\mathrm{d}}
\begin{document}
\section*{UnsolvedMath ID: 30000742}
\textbf{Source problem code:} catalogue entry
\subsection{Definitions and mathematical statement}
For a graph on a prescribed vertex set, let $H$ be the property of containing a Hamilton cycle, let $D_k$ be the property that every vertex has degree at least $k$, and let $C_k$ denote $k$-connectivity. Under the usual vertex-connectivity convention, $C_k$ means that the graph has more than $k$ vertices and remains connected after deletion of any set of fewer than $k$ vertices. The source does not specify whether this convention or edge connectivity is intended. Write $b_P$ for the critical bias of the positional game with target property $P$, using one fixed board and one fixed bias/turn convention for all compared properties. The two source questions are \[b_H=b_{D_2}\ ?\qquad\text{and}\qquad b_{C_k}=b_{D_k}\ ?\] These are written as equalities of critical biases, not merely equality of asymptotic leading terms. For example, in the customary $(1:b)$ Maker--Breaker edge game on $K_n$, Maker claims one free edge and Breaker claims $b$ free edges per round, and Maker seeks to achieve $P$ in Maker's own spanning graph; a common convention sets $b_P(n)$ to the largest integer $b$ for which Maker has a winning strategy, with zero if there is no positive winning bias. Under that convention the displayed questions compare $b_H(n)$ with $b_{D_2}(n)$ and $b_{C_k}(n)$ with $b_{D_k}(n)$. This is an explicit conditional illustration: the supplied extract does not identify the board, game type, first player, critical-bias convention, or quantifier range in $n$ and $k$, so it does not determine whether the intended equalities hold for every admissible size or only eventually. The supplied literature triage of 2026-08-21 labels the questions open and reports no definitive resolution located.
\subsection{Short English statement}
Ask whether the Hamiltonicity threshold bias equals the minimum-degree-two threshold, and whether the $k$-connectivity threshold equals the minimum-degree-$k$ threshold. The exact game conventions and parameter ranges are missing from the extract.
\subsection{Sources}
\begin{itemize}\item Local UnsolvedMath registry, permanent ID 30000742, source problem catalogue entry, \emph{Equalities Among Critical Biases in Positional Games}, proposed year 2007; no individual proposer supplied.\item Supplied citation: \emph{Mini-Workshop: Positional Games} (2008), Oberwolfach Reports, DOI 10.4171/owr/2007/20, pp.\ 1073--1114: \url{https://doi.org/10.4171/owr/2007/20}. Supplied literature label: original Oberwolfach report.\end{itemize}
\end{document}
